\lecture{7}
\title{Decision Problems for Automata and Grammars}

\begin{document}

By the end of this lecture, we will have gone over
all content necessary for PSET \#2 and our quiz next week.

\subsection{Encodings}

\textbf{Definition.}
Given \emph{any} object $z$
(e.g., an automaton, a graph, a string),
let $\langle z \rangle$ denote an \textbf{encoding} of $z$
into $\{\str{0}, \str{1}\}^*$.
Furthermore, let $\langle z_1, z_2, \dots, z_k \rangle$
denote an encoding of
$(\langle z_1 \rangle, \langle z_2 \rangle, \dots, \langle z_k \rangle)$
into a single string.

\begin{remark}
    Note that $\langle z_1, z_2 \rangle$ should probably not just be
    the concatenation of $\langle z_1 \rangle$ and $\langle z_2 \rangle$.
    The encoding must be such that one can read the binary string
    $\langle z_1, z_2 \rangle$ and recognize that
    it must be an encoding of two objects,
    not just $\langle z_3 \rangle$ for some random $z_3$.
\end{remark}

The details of the encoding are not important---their
only purpose is to let us take \emph{anything}
as binary input to a TM.

\subsection{Algorithms on Automata}

Turing machines are really strong---they
can even run algorithms that answer questions about algorithms!

\begin{problem}
    Consider the language
    $A_{\mathrm{DFA}} := \{\langle B, w \rangle \mid
    B \text{ is a DFA that accepts } w\}$.
    Show that $A_{\mathrm{DFA}}$ is decidable.
\end{problem}

\begin{solution}
    Recall that it suffices to define a TM decider $M$
    for which $A_{\mathrm{DFA}} = L(M)$.
    Here's what $M$ does on input $s$.
    \begin{enumerate}
        \item Check that $s$ is of the form $\langle B, w \rangle$
        for some DFA $B$ and some string $w$.
        \item Simulate $B$ on $w$.
        (This takes finite time,
        because $B$ is a \emph{DFA} that takes
        a \emph{finite} string $w$ as input.)
        \item In the end, return \textsc{Accept}
        if $B$ returns \textsc{Accept},
        and return \textsc{Reject} otherwise.
    \end{enumerate}
    And this works, obviously. ~ $\blacksquare$
\end{solution}

We're \emph{really} abstracting away the details here.
But because of the Church-Turing thesis---an algorithm
\emph{is} a Turing machine---we should have confidence
that working at such a high level of abstraction
will not cause any issues.

\begin{remark}
    In case you're not convinced,
    here's a sketch of how we might implement
    ``simulate $B$ on $w$'' in a TM.
    \begin{enumerate}
        \item Use three tapes to hold each of
        $\langle B, w \rangle$, the current state $q \in Q$,
        and $w$ with its head at the current position.
        \item At each step, do the following:
        \begin{enumerate}
            \item Read the symbol $a \in \Sigma$ at tape 3.
            \item Scan the transition table on tape 1 for the entry
            whose state matches tape 2 and whose symbol is $a$.
            \item Write the target state onto tape 2
            and advance the tape 3 head.
        \end{enumerate}
        \item When all of $w$ has been read,
        scan $F$ on tape 1 for the state on tape 2.
    \end{enumerate}
    Perhaps the individual steps above
    are more believably implementable by a TM.
\end{remark}

\begin{problem}
    Consider the language
    $A_{\mathrm{NFA}} := \{\langle B, w \rangle \mid
    B \text{ is an NFA that accepts } w\}$.
    Show that $A_{\mathrm{NFA}}$ is decidable.
\end{problem}

\begin{solution}
    Just take a TM decider $M$ that, on input $s$,
    does the following:
    \begin{quote}
        Check that $s$ takes form $\langle B, w \rangle$,
        then convert the NFA $B$ to a DFA $B'$,
        then simulate $B'$ on $w$ as in the previous problem.
        ~ $\blacksquare$
    \end{quote}
\end{solution}

\begin{problem}
    Consider the language
    $E_{\mathrm{DFA}} := \{\langle B \rangle \mid
    B \text{ is a DFA and } L(B) = \emptyset\}$.
    Show that $E_{\mathrm{DFA}}$ is decidable.
\end{problem}

\begin{solution}
    This is a path-finding problem---it asks the following:
    \begin{quote}
        \textbf{Question.}
        Is there a path from the start state to any accept state?
    \end{quote}
    So we can just return \textsc{Accept} if and only if
    BFS from the start state finds \emph{no} accept state. ~ $\blacksquare$
\end{solution}

\begin{problem}
    Consider the language $EQ_{\mathrm{DFA}} := \{
        \langle A, B \rangle \mid
        A \text{ and } B \text{ are DFAs with }
        L(A) = L(B)
    \}$.  Show that $EQ_{\mathrm{DFA}}$ is decidable.
\end{problem}

\begin{solution}
    This one is less obvious.
    Recall from our ``closure property constructions'' from
    \href{/18.404/lectures/1/#regular-operations-and-regular-expressions}{Lecture 1}
    that
    \[ L(C) := \left(
        L(A) \cap \overline{L(B)}
    \right) \cup \left(
        \overline{L(A)} \cap L(B)
    \right) \]
    for some DFA $C$,
    where $C$ is \emph{constructible} from the DFAs $A$ and $B$.

    Then we just consider a TM that constructs the DFA $C$,
    then runs the algorithm used for $E_{\mathrm{DFA}}$.
    ~ $\blacksquare$
\end{solution}

\subsection{Algorithms on Grammars and Turing Machines}

Answering questions about grammars is a little harder.
Here's a useful fact that we'll assume without proof.

\begin{theorem}{Chomsky Normal Form}
    Any CFG can be expressed as a CFG
    in which all rules look like $A \rightarrow BC$ or $A \rightarrow a$.
    Furthermore, there is a TM-executable algorithm
    for converting any CFG into such a form.
\end{theorem}

\begin{problem}
    Consider the language
    $A_{\mathrm{CFG}} := \{\langle G, w \rangle \mid
    G \text{ is a CFG that generates } w\}$.
    Show that $A_{\mathrm{CFG}}$ is decidable.
\end{problem}

\begin{solution}
    First check that input $s$ takes form $\langle G, w \rangle$.
    Then convert $G$ into Chomsky Normal Form.

    The point is that any derivation of $w$ must apply a rule
    exactly $2|w| - 1$ times.  So it is now feasible
    to search through all derivations of length $2|w| - 1$
    and \textsc{Accept} if and only if any derivation yields $w$.
    ~ $\blacksquare$
\end{solution}

Here's an interesting corollary to the above point.

\begin{theorem}{CFLs are Decidable}
    Every context-free language is decidable.
\end{theorem}

\begin{proof}
    We must show that if $A$ is a CFL
    generated by a CFG $G$, then
    there is a TM $M_G$ that decides $A$.
    Well, just take $M_G$ to be the TM that,
    on input $w$, searches through the
    length-$(2|w| - 1)$ derivations using the
    Chomsky Normal Form of $G$.
    ~ $\blacksquare$
\end{proof}

\begin{problem}
    Consider the language
    $E_{\mathrm{CFG}} := \{\langle G \rangle \mid
    G \text{ is a CFG and } L(G) = \emptyset\}$.
    Show that $E_{\mathrm{CFG}}$ is decidable. \\
    For example, the CFG $G$ with only one rule
    $S \rightarrow aS$ should be in $E_{\mathrm{CFG}}$.
\end{problem}

\begin{solution}
    Construct the CFG $G'$ that results from
    replacing every terminal in $G$ with $\epsilon$.
    (So $S \rightarrow \str{0}S\str{1}$
    becomes $S \rightarrow \epsilon S \epsilon$, say.)
    Then just run the algorithm used for $A_{\mathrm{CFG}}$
    to check whether $G'$ generates $\epsilon$,
    and \textsc{Accept} if and only if it does not. ~ $\blacksquare$

    \dots except, the algorithm for $A_{\mathrm{CFG}}$
    now must search through
    all derivations of length $2|\epsilon| - 1 = -1$.
    What does that mean?!

    The reason for this confusion is that
    Chomsky Normal Form permits one more kind of rule:
    $S \to \epsilon$, where $S$ is the start variable.
    So all the algorithm $A_{\mathrm{CFG}}$ is doing
    is determining whether $S \rightarrow \epsilon$
    is in the Chomsky Normal Form of $G'$.

    And that's really unsatisfying---all the nontrivial work
    is offloaded to the conversion of $G'$ into Chomsky Normal Form.
    But we don't have time to present a better proof; maybe next lecture.
    ~ $\blacksquare$
\end{solution}

It turns out that, unlike with DFAs, the following languages
are \emph{not} decidable.
\begin{itemize}
    \item $EQ_{\mathrm{CFG}} := \{\langle G, H \rangle \mid
    G \text{ and } H \text{ are CFGs with }
    L(G) = L(H)\}$.
    \item $A_{\mathrm{TM}} := \{\langle M, w \rangle \mid
    M \text{ is a TM that accepts } w\}$.
\end{itemize}
But if we replace ``decidable'' with ``Turing-recognizable'',
we find\dots

\begin{problem}
    Show that $A_{\mathrm{TM}}$ is Turing-recognizable.
\end{problem}

\begin{solution}
    After checking that input $s$ takes form $\langle M, w \rangle$,
    just simulate $M$ on $w$.  In particular,
    \begin{itemize}
        \item If $M$ halts and accepts, then return \textsc{Accept}.
        \item If $M$ halts and rejects, then return \textsc{Reject}.
        \item If $M$ never halts\dots well, then
        this Turing machine won't halt, either.
    \end{itemize}
    This just works. ~ $\blacksquare$
\end{solution}

\end{document}